\documentclass[12pt,a4paper]{article}

% --- Packages ---
\usepackage[utf8]{inputenc}
\usepackage[margin=1in]{geometry}
\usepackage{graphicx}
\usepackage{array}
\usepackage{hyperref}
\usepackage{setspace}
\usepackage{xcolor}

% --- Document Settings ---
\onehalfspacing
\hypersetup{colorlinks=true, linkcolor=blue, urlcolor=blue, citecolor=black}

% --- Title Page ---
\begin{document}

\begin{titlepage}
    \centering
    \vspace*{1cm}
    {\Large \textbf{Research Project Proposal (Milestone 1)}}\\[1.5cm]

    
    {\huge \textbf{District-Level Dengue Outbreak Early Warning System for Sri Lanka.}}\\[3cm]
    
    \textbf{Group Number: Group 10}\\[5cm]
    
    {\textbf{Group Members:}}\\[0.3cm]
    \begin{table}[h]
        \centering
        \begin{tabular}{|l|l|l|}
        \hline
        \textbf{Name} & \textbf{Student ID} \\ \hline
        Kavindu Amalka & SE/2021/008 \\ \hline
        Hiruni Imasha & SE/2021/053 \\ \hline
        Sandaru Mihiran & SE/2021/010 \\ \hline
        Sachini Weerakkody & SE/2021/018 \\ \hline
        \end{tabular}
    \end{table}
    
    \vfill
    Date: \today
\end{titlepage}



% --- 1. INTRODUCTION ---
\section{Introduction}

\subsection{Background}
\textit{Describe the basic facts and provide context about the research area.}

[Provide context about dengue in Sri Lanka. Include:
\begin{itemize}
    \item Current prevalence and impact of dengue
    \item Disease characteristics and transmission
    \item Geographic distribution in Sri Lanka
    \item Seasonal patterns and risk factors
\end{itemize}
]

\subsection{Importance}
\textit{Highlight the significance of the research for industry practices or advancing knowledge.}

[Explain why this research matters:
\begin{itemize}
    \item Public health impact and burden of disease
    \item Current limitations in outbreak prediction
    \item Potential benefits of early warning systems
    \item Industry or societal need
\end{itemize}
]

\subsection{Motivation}
\textit{Explain why this research area is of interest to you and its potential contributions.}

[Describe personal/team motivation:
\begin{itemize}
    \item Why your team chose this topic
    \item Potential contributions to the field
    \item Practical solutions you aim to provide
\end{itemize}
]

% --- 2. PROBLEM STATEMENT ---
\section{Problem Statement}

\subsection{The Core Challenge: Reactive vs. Proactive}
Dengue management in Sri Lanka is largely reactive, triggering interventions only after infection rates spike. While countries like Singapore and Brazil use Machine Learning (ML) to anticipate outbreaks, Sri Lanka lacks a localized, proactive early warning system that provides the necessary lead time for effective prevention.

\subsection{Key Gaps and Challenges}
\begin{itemize}
    \item \textbf{Data Fragmentation:} Existing models rely primarily on rainfall, failing to integrate multidimensional factors like Human Mobility and PHI Larval Surveillance data for high-accuracy predictions.
    \item \textbf{Geographical Oversimplification:} Literature often overlooks the distinct dengue-climate dynamics between Wet-Zone and Dry-Zone districts and the separate impacts of the Southwest and Northeast monsoons.
    \item \textbf{Lead-Time Inefficiency:} Without a district-specific lag analysis, Public Health Inspectors (PHIs) cannot time their vector control interventions precisely, resulting in inefficient resource allocation.
    \item \textbf{Technical Limitations:} Traditional models (e.g., ARIMA) struggle with the non-linear temporal patterns of long-term data. Advanced architectures like LSTM and XGBoost are essential to model these complex 7-year longitudinal relationships.
\end{itemize}

\subsection{Proposed Solution: Actionable Prevention}
This research introduces a Dual-Stream Intelligence framework integrating human mobility and PHI larval reports with micro-climatic data. By utilizing LSTM and XGBoost architectures, the system converts complex epidemiological patterns into a 4-week-ahead predictive dashboard. This tool transforms raw data into actionable insights, empowering health authorities to transition from crisis management to proactive strategic prevention.

% --- 3. OBJECTIVES ---
\section{Objectives}

\textit{Outline the specific goals your research aims to achieve.}

\subsection{Primary Objective}
[State your main research goal clearly and measurably]

\subsection{Secondary Objectives}
\begin{enumerate}
    \item Objective 1: Specific, measurable goal
    \item Objective 2: Specific, measurable goal
    \item Objective 3: Specific, measurable goal
    \item Objective 4: Specific, measurable goal
\end{enumerate}

\noindent\textbf{Guidelines:}
\begin{itemize}
    \item Explain the benefits or impacts of your research (e.g., improved understanding, enhanced performance, practical solutions)
    \item Ensure objectives are clear, measurable, and directly address the problem statement
\end{itemize}

% --- 4. PRELIMINARY LITERATURE REVIEW ---
\section{Preliminary Literature Review}

\subsection{Existing Research and Key Findings}
Dengue forecasting in Sri Lanka has been a significant area of study, with research consistently linking disease outbreaks to climatic factors such as rainfall, temperature, and humidity \cite{wickramathilaka2024}. These environmental variables directly influence the breeding cycles and population density of the Aedes mosquito \cite{withanage2018}.

\begin{itemize}
    \item \textbf{Methodologies:} Recent research has shifted toward Deep Learning techniques, particularly Long Short-Term Memory (LSTM) networks. These models are favored for their ability to handle non-linear patterns and long-term dependencies in historical data \cite{albannoud2026, uduwanage2025}.
    \item \textbf{Strengths:} Current forecasting models effectively demonstrate the correlation between weather fluctuations and immediate spikes in dengue cases, providing a foundation for early warning systems \cite{karasinghe2024}.
    \item \textbf{Weaknesses:} However, much of the existing research remains theoretical and focused on model accuracy rather than practical use. There is a notable lack of functional software dashboards that health authorities can use in real-time \cite{wickramathilaka2024}. Additionally, many studies focus primarily on the Western Province, which may not accurately reflect the disease dynamics in other ecological zones of the country \cite{fernando2024}.
\end{itemize}

\subsection{Research Gap and Justification}
A review of prior literature highlights a critical gap in turning high-accuracy predictive models into actionable, software-driven tools for the public health sector \cite{wickramathilaka2024}.

\begin{itemize}
    \item \textbf{Addressing the Gaps:} This research addresses these limitations by developing a \textbf{Functional Predictive Dashboard}. To move beyond the common "Western Province Bias," our study analyzes \textbf{one critical district per province}, ensuring the tool is relevant and accurate for the entire country \cite{fernando2024}.
    \item \textbf{Justification:} This study is justified by the urgent need for proactive forecasts that assist in efficient resource allocation during outbreaks \cite{rahulan2024}. By utilizing a weekly panel dataset since 2017, we diverge from climate-only models by also incorporating human mobility and PHI larval reports, which are crucial but often overlooked factors in dengue transmission \cite{delaguila2024}.
\end{itemize}

% --- 5. METHODOLOGY ---
\section{Methodology}

\subsection{Research Approach}
\textit{Define the research methods and logical steps you will take to address the problem.}

[Describe your overall approach:
\begin{itemize}
    \item Research design (e.g., experimental, case study, modeling, survey)
    \item Why this approach is appropriate
    \item Justification for methodological choices
\end{itemize}
]

\subsection{Data Collection and Analysis}
\textit{Detail the approach you will use to collect and analyze data.}

[Specify:
\begin{itemize}
    \item Data sources and collection methods
    \item Data types and characteristics
    \item Data analysis techniques
    \item Tools and software to be used (e.g., Python, TensorFlow, R)
\end{itemize}
]

\subsection{Plan for Achieving Research Objectives}
\textit{Provide a clear plan for achieving the research objectives.}

[Detail:
\begin{itemize}
    \item Step-by-step approach aligned with objectives
    \item Tools, technologies, and techniques
    \item Expected outputs and deliverables for each phase
\end{itemize}
]

% --- 6. PROJECT SCHEDULE ---
\section{Project Schedule}

\textit{Develop a detailed project timeline with key milestones and assign team members.}

\subsection{Gantt Chart}
The project timeline from initiation to completion is visualized below:

\begin{table}[h]
\centering
\resizebox{\textwidth}{!}{%
\begin{tabular}{|l|c|c|c|c|c|c|}
\hline
\textbf{Milestone/Phase} & \textbf{M1} & \textbf{M2} & \textbf{M3} & \textbf{M4} & \textbf{M5} & \textbf{Member(s) In-charge} \\ \hline
[Phase 1: Literature Review] & X &  &  &  &  & [Member Names] \\ \hline
[Phase 2: Data Collection/Prep] &  & X &  &  &  & [Member Names] \\ \hline
[Phase 3: Model Development] &  &  & X &  &  & [Member Names] \\ \hline
[Phase 4: Testing \& Validation] &  &  &  & X &  & [Member Names] \\ \hline
[Phase 5: Final Documentation] &  &  &  &  & X & [Member Names] \\ \hline
\end{tabular}%
}
\caption{Project Timeline and Milestones}
\end{table}

\subsection{Milestone Definitions}
\begin{itemize}
    \item \textbf{[Milestone 1]:} [Description of measurable outcome]
    \item \textbf{[Milestone 2]:} [Description of measurable outcome]
    \item \textbf{[Milestone 3]:} [Description of measurable outcome]
    \item \textbf{[Milestone 4]:} [Description of measurable outcome]
    \item \textbf{[Milestone 5]:} [Description of measurable outcome]
\end{itemize}

\noindent\textbf{Guidelines:}
\begin{itemize}
    \item Identify key milestones covering all major phases from initiation to completion
    \item Assign team members to each milestone
    \item Ensure milestones are measurable and achievable, representing significant outcomes or deliverables
\end{itemize}

% --- 7. REFERENCES ---
\section{References}

\noindent\textit{Include all sources cited in your proposal using APA citation style. Minimum 10 recent scholarly sources required.}

\begin{enumerate}
    % 1. Machine Learning Forecasting (Brazil Study)
    \item Al Bannoud, M., Mendes, J. L., Sacay, R. S., Martins, T. D., Bresolin, I. R. A. P., & Lourenço, J. B. (2026). Machine Learning forecasting of dengue in São Paulo using virtual data augmentation and urban incident predictors. \textit{Acta Tropica}, \textit{275}, 107992. https://doi.org/10.1016/j.actatropica.2024.107992
    
    % 2. System Dynamics & Risk Mapping
    \item Del-Águila-Mejía, J., Morilla, F., & Donado-Campos, J. D. M. (2024). A system dynamics modelling and analytical framework for imported dengue outbreak surveillance and risk mapping. \textit{Acta Tropica}, \textit{257}, 107304. https://doi.org/10.1016/j.actatropica.2024.107304
    
    % 3. Systematic Review (Sri Lanka 2000-2020)
    \item Fernando, L., Wickramasinghe, V. P., Kalra, P., Kastner, R., & Gallagher, E. (2024). Epidemiological burden of dengue in Sri Lanka: A systematic review of literature from 2000-2020. \textit{IJID Regions}, \textit{13}, 100436. https://doi.org/10.1016/j.ijregi.2024.100436
    
    % 4. Modified ARIMA Modeling (Sri Lanka)
    \item Karasinghe, N., Peiris, S., Jayathilaka, R., & Dharmasena, T. (2024). Forecasting weekly dengue incidence in Sri Lanka: Modified Autoregressive Integrated Moving Average modeling approach. \textit{PLOS ONE}, \textit{19}(3), e0299953. https://doi.org/10.1371/journal.pone.0299953
    
    % 5. Case Study (Gampaha District)
    \item Rahulan, K., Ranagala, R. M., & Paul Roshan, G. (2024). Dengue Outbreak Control and Prevention: A Case Study of Gampaha District, Sri Lanka. \textit{The Sri Lanka Journal of Medical Administration}, \textit{25}(1), 12-21. https://doi.org/10.4038/sljma.v25i1.5477
    
    % 6. ML Prediction (University of Peradeniya Study)
    \item Uduwanage, H. U., Konara, K. M. S. L., Mihiranga, G. D. R., Karunarathna, S. N., Noordeen, F., & Ekanayake, I. (2025). Prediction of dengue outbreaks in Sri Lanka using machine learning techniques. \textit{Sri Lanka Journal of Medicine}, \textit{34}(1), 15-26. https://doi.org/10.4038/sljm.v34i1.568
    
    % 7. White Paper (University of Moratuwa)
    \item Wickramathilaka, S., & Munasinghe, A. (2024). \textit{Epidemiological pattern recognition and disease forecasting for dengue in Sri Lanka} [White paper]. University of Moratuwa.
    
    % 8. Gampaha Forecasting Model (Basic Model)
    \item Withanage, G. P., Viswakula, S. D., Silva Gunawardena, Y. I. N., & Hapugoda, M. D. (2018). A forecasting model for dengue incidence in the District of Gampaha, Sri Lanka. \textit{Parasites \& Vectors}, \textit{11}(1), 262. https://doi.org/10.1186/s13071-018-2828-2
    
    % 9. Interpretable Modeling (Singapore Study)
    \item Zhang, Z., Hu, S., Hu, J., Hao, Z., & Huang, J. (2026). Decoding multi-frequency scale climate-dengue dynamics: An interpretable modeling framework. \textit{Urban Climate}, \textit{65}, 102734. https://doi.org/10.1016/j.uclim.2024.102734
    
    % 10. One Health Concept (Dengue Control)
    \item Feng, X., Jiang, N., Zheng, J., Zhu, Z., Chen, J., & Duan, L. (2024). Advancing knowledge of One Health in China: Lessons for One Health from China's dengue control and prevention programs. \textit{Science in One Health}, \textit{3}, 100087. https://doi.org/10.1016/j.soh.2024.100087
\end{enumerate}

\end{document}